\documentclass[11pt,a4paper]{article}
\usepackage[margin=22mm]{geometry}
\usepackage{amsmath,amssymb}
\usepackage{fontspec}
\usepackage{polyglossia}
\setdefaultlanguage{russian}
\setmainfont{DejaVu Serif}
\usepackage[unicode,hidelinks]{hyperref}

\allowdisplaybreaks[1]
\raggedright
\setlength{\parskip}{0.35em}
\newcommand{\dd}{\mathop{}\!\mathrm{d}}
\newcommand{\avg}[1]{\langle #1\rangle}
\newcommand{\ket}[1]{\lvert #1\rangle}
\newcommand{\bra}[1]{\langle #1\rvert}
\newcommand{\braket}[2]{\langle #1\mid #2\rangle}
\DeclareMathOperator{\Var}{Var}

\newcommand{\solutiontitle}[4]{%
  \hypersetup{pdftitle={Квантовая механика 1. Решения ДЗ №#1}}%
  \begin{center}
    {\small Квантовая механика 1 \enspace\textperiodcentered\enspace Семестр 1\par}
    \vspace{1ex}
    {\LARGE\bfseries Решения домашнего задания №#1\par}
    \vspace{0.5ex}
    {\large #2\par}
  \end{center}
  \noindent{\small Условия приведены в файле «#3», #4.}\par\medskip
}

\begin{document}
\solutiontitle{9}{Кусочные и периодические потенциалы}{ДЗ2026(1)\_09 (1).pdf}{с.~1}

\section{Дельта-яма в центре бесконечной ямы}

Пусть \(L=a/2\), \(g=m\alpha/\hbar^2>0\).
Граничные условия \(\Psi(0)=\Psi(a)=0\). Через центр функция
непрерывна, а из интегрирования уравнения Шрёдингера следует
\[
 \Psi'(L^+)-\Psi'(L^-)=-2g\Psi(L).
\]
Потенциал симметричен относительно \(L\), поэтому можно решить задачу
по отдельности для чётности \(s=\pm1\) относительно \(L\).
При \(E>0\) положим \(k=\sqrt{2mE}/\hbar\). Левая половина
\(\Psi(x)=A\sin(kx)\), правая \(\Psi(x)=sA\sin(k(a-x))\).

Для нечётных состояний непрерывность требует \(\Psi(L)=0\).
Производная тогда непрерывна автоматически, так что \(kL=n\pi\),
\(n=1,2,\ldots\). Их энергии \(\hbar^2(2n\pi/a)^2/(2m)\).
Для чётных состояний скачок равен \(-2Ak\cos(kL)\).
Следовательно \(k\cos(kL)=g\sin(kL)\). Одновременный нуль
синуса и косинуса невозможен, и уравнение принимает вид
\[
 z\cot z=gL,\qquad z=kL,\qquad E=\frac{\hbar^2 z^2}{2mL^2}.
\]
На интервалах \((n\pi,n\pi+\pi/2)\) при \(n\ge1\) левая часть
убывает от бесконечности до нуля, поэтому на каждом есть один корень.
Убывание следует из производной
\((\sin z\cos z-z)/\sin^2z<0\).
На \((0,\pi/2)\) диапазон значений равен \((0,1)\), и
положительный корень существует только при \(gL<1\).

При \(E<0\) берём \(\kappa=\sqrt{-2mE}/\hbar\) и на левой
половине \(A\sinh(\kappa x)\). У неё нет нуля при \(x=L>0\),
поэтому возможна только чётная ветвь. Скачок даёт
\[
 \kappa\coth(\kappa L)=g,\qquad E=-\frac{\hbar^2\kappa^2}{2m}.
\]
Функция \(u\coth u\) возрастает от \(1\) до бесконечности
при \(u>0\), поскольку \(\sinh u\cosh u>u\).
Итак, отрицательный уровень ровно один при \(gL>1\).
При \(E=0\) левая половина линейна, \(Ax\).
Чётное отражение даёт \(-2A=-2gAL\), то есть нулевой уровень
возникает при \(gL=1\). У нечётного отражения непрерывность
потребовала бы \(AL=0\). Для нормировки всех найденных функций
достаточно взять \(|A|^{-2}=2\int_0^L f(x)^2\dd x\),
где \(f(x)=\sin(kx),\sinh(\kappa x)\) или \(x\) соответственно.

\section{Бесконечная яма со ступенькой}

Стенки задают \(\Psi(0)=\Psi(a)=0\). При \(U_0>0\)
энергия \(E>0\), потому что
\(E=\hbar^2\int|\Psi'|^2\dd x/(2m)+\int U|\Psi|^2\dd x\)
для нормированного состояния, а обе части неотрицательны.
Первая строго положительна для ненулевой функции, обращающейся
в ноль на стенках. Обозначим \(L=a/2,\ k=\sqrt{2mE}/\hbar\).
На \(x<L\) функция \(A\sin(kx)\). На \(x>L\) используем
решение, обнуляющееся в \(a\), с множителем \(B\).

При \(E>U_0\) оно равно \(B\sin(k_2(a-x))\), где
\(k_2=\sqrt{2m(E-U_0)}/\hbar\). Сшивка значения и производной даёт
\[
 A\sin(kL)=B\sin(k_2L),\qquad Ak\cos(kL)=-Bk_2\cos(k_2L).
\]
Приравнивая определитель этой системы к нулю, получаем
\[
 k\cos(kL)\sin(k_2L)+k_2\sin(kL)\cos(k_2L)=0.
\]
Уравнение сохраняет уровни с узлом на ступеньке, когда оба синуса
равны нулю и отношение \(B/A\) задаёт вторая строка сшивки.

При \(0<E<U_0\) правая функция \(B\sinh(\kappa(a-x))\),
\(\kappa=\sqrt{2m(U_0-E)}/\hbar\).
Сшивка даёт \(A\sin(kL)=B\sinh(\kappa L)\) и
\(Ak\cos(kL)=-B\kappa\cosh(\kappa L)\).
Нуль левого синуса привёл бы к \(B=0\) и далее \(A=0\).
Поэтому уравнение на энергию имеет вид
\[
 k\cot(kL)=-\kappa\coth(\kappa L).
\]
Положим \(z=kL,\ z_0=L\sqrt{2mU_0}/\hbar,\ s=\sqrt{z_0^2-z^2}\).
Уравнение принимает вид \(-z\cot z=s\coth s\).
На ветви \(((j+\tfrac12)\pi,(j+1)\pi)\) левая часть
возрастает от нуля до бесконечности. Правая часть убывает
с ростом \(z\) и стремится к \(1\) при \(z=z_0\).
Поэтому на ветви есть ровно один корень с \(E<U_0\), если
\(z_0\) больше точки этой ветви, в которой \(-z\cot z=1\).
При \(E=U_0\) справа \(B(a-x)\) и уравнение переходит в
\(k\cot(kL)=-1/L\), в согласии с пределом \(s\to0\).
В каждом случае корень даёт состояние на конечном интервале,
а отношение \(B/A\) фиксируют уравнения сшивки.
Последнюю константу задаёт
\(|A|^2 I_L+|B|^2 I_R=1\).
Интегралы квадратов левой и правой функций на соответствующих
половинах вычисляются как \(L/2-\sin(2kL)/(4k)\) для синуса
с волновым числом \(k\), \(\sinh(2\kappa L)/(4\kappa)-L/2\)
для гиперболического синуса и \(L^3/3\) для линейной функции.

\section{Модель Кронига --- Пенни}

Для участка постоянного потенциала с волновым числом \(w\)
и длиной \(l\) матрица перехода для
\((\Psi,\Psi')^{\mathsf T}\) следует из решения \(\Psi''=-w^2\Psi\).
\[
 M(w,l)=
 \begin{pmatrix}
 \cos(wl)&\sin(wl)/w\\ -w\sin(wl)&\cos(wl)
 \end{pmatrix},\qquad \det M=1.
\]
При \(w=0\) берём непрерывный предел, а при мнимом \(w\)
тригонометрические функции переходят в гиперболические.
Положим \(d=b-a,\ w_1^2=2m(E+U_0)/\hbar^2,\ w_2^2=2mE/\hbar^2\).
Сшивка без скачка производной даёт матрицу периода
\(M_b=M(w_2,d)M(w_1,a)\). По теореме Блоха её собственное
значение равно \(e^{iqb/\hbar}\), где \(q\) имеет размерность
импульса. Второе собственное значение равно \(e^{-iqb/\hbar}\),
поскольку \(\det M_b=1\). Следовательно
\begin{align*}
 \cos(qb/\hbar)=F(E)
 &=\tfrac12\operatorname{tr}M_b\\
 &=\cos(w_1a)\cos(w_2d)
   -\frac{w_1^2+w_2^2}{2}\,
     \frac{\sin(w_1a)}{w_1}\frac{\sin(w_2d)}{w_2}.
\end{align*}
При \(-U_0<E<0\), \(w_2=i\kappa\), второй косинус заменяется
на \(\cosh(\kappa d)\), а \(\sin(w_2d)/w_2\) на
\(\sinh(\kappa d)/\kappa\).
Матрица периода вещественна для вещественного \(E\).
Условие \(|F(E)|\le1\) эквивалентно наличию собственного значения
на единичной окружности. Каждый его собственный вектор задаёт
начальные данные для ненулевой блоховской волны, в том числе
при \(F(E)=\pm1\). Это разрешённые зоны. Нижняя граница энергии
не может быть ниже \(-U_0\) по положительности кинетической формы
на ячейке с блоховскими граничными условиями.

\textbf{Численный пример.}
Параметры в листке не фиксированы, поэтому берём \(a/b=1/2\) и
\(U_0=8E_b\), \(E_b=\hbar^2/(2mb^2)\).
Положим \(\epsilon=E/E_b,\ Q=qb/\hbar\). В формуле для \(F\)
можно измерять длины в единицах \(b\), с обоими отрезками
длины \(1/2\) и квадратами волновых чисел \(\epsilon+8,\epsilon\).
Ищем первые три корня \(F(E_b\epsilon)-\cos Q=0\) по возрастанию
\(\epsilon\). Для каждой строки таблицы достаточно просканировать
\([-8,85]\) с шагом \(10^{-3}\), выделить интервалы смены знака и
уточнить корни методом Брента. Отрицательные \(q\) дают те же
числа, поскольку правая часть зависит от \(\cos Q\).
\[
 \begin{array}{c|rrr}
 Q/\pi & E_1/E_b & E_2/E_b & E_3/E_b\\ \hline
 0   &-4.33073 &35.37717 &35.77975\\
 1/4 &-3.73512 &26.37342 &46.03683\\
 1/2 &-1.96832 &18.45548 &57.73390\\
 3/4 & 0.83749 &11.95745 &70.66302\\
 1   & 3.23976 & 8.32423 &84.02171
 \end{array}
\]
Для проверки порядка корней диагонализируем ту же задачу
на ячейке в базисе \(e^{i(2\pi n+Q)y}\), \(y=x/b\).
В единицах \(E_b\) матрица равна
\[
 H_{nm}=(2\pi n+Q)^2\delta_{nm}
            -8\int_0^{1/2}e^{2\pi i(m-n)y}\,\dd y .
\]
При \(|n|,|m|\le110\) и \(220\) три нижних собственных значения
совпадают с таблицей с точностью до округления.

\section{Дираковская гребёнка}

Пусть \(\alpha>0,\ g=m\alpha/\hbar^2\). Через каждый
барьер \(\Psi\) непрерывна, а \(\Psi'\) увеличивается на \(2g\Psi\).
Значит матрица скачка \(J=\begin{pmatrix}1&0\\2g&1\end{pmatrix}\),
а матрица от точки сразу справа от одного барьера до точки сразу
справа от следующего равна \(JM(k,a)\), \(k=\sqrt{2mE}/\hbar\).
По её следу и определителю получаем
\[
 \cos(qa/\hbar)=\cos z+\frac{ga}{z}\sin z,\qquad z=ka.
\]
Разрешены энергии, при которых модуль правой части не превосходит
единицы. При \(E=0\) предел равен \(1+ga>1\),
при \(E<0\) получаем \(\cosh(\kappa a)+(g/\kappa)\sinh(\kappa a)>1\).
При \(z=n\pi+u\), \(n=0,1,\ldots,\ 0<u<\pi\),
умножаем правую часть на \((-1)^n\) и получаем
\(\cos u+(ga/z)\sin u>-1\).
Условие сверху \(\cos u+(ga/z)\sin u\le1\) после деления на
\(\sin u>0\) становится \(\tan(u/2)\ge ga/(n\pi+u)\).
Тангенс строго возрастает от нуля до бесконечности, правая
часть убывает, так что уравнение равенства имеет один корень \(u_n\).
Учитывая правые края, где \(\sin z=0\), получаем зоны
\[
 \frac{\hbar^2(n\pi+u_n)^2}{2ma^2}\ \le E\le\
 \frac{\hbar^2(n+1)^2\pi^2}{2ma^2},\qquad n=0,1,\ldots .
\]
В пределе \(\alpha\to0^+\) имеем \(u_n\to0\) и щели закрываются.

\section{Свободное расплывание двух гауссовых пакетов}

Пусть \(C=[2\sigma\sqrt\pi(1+e^{-x_0^2/\sigma^2})]^{-1/2}\)
и \(\sigma>0\). Интегралы квадратов двух начальных гауссиан
равны \(\sigma\sqrt\pi\), а их перекрытие
\(\sigma\sqrt\pi e^{-x_0^2/\sigma^2}\), так что \(C\)
соответствует норме один. В импульсном пространстве свободная
эволюция умножает каждую гауссиану на \(e^{-ip^2t/(2m\hbar)}\).
Вместе с начальным гауссовым множителем это даёт
\(\exp[-\sigma^2p^2 D/(2\hbar^2)]\), где
\(\tau=\hbar t/(m\sigma^2)\), \(D=1+i\tau\).
Обратный гауссов интеграл с фазами сдвига возвращает
\[
 \Psi(x,t)=\frac C{\sqrt D}(g_++g_-),\qquad
 g_\pm=\exp\!\left[-\frac{(x\mp x_0)^2}{2\sigma^2 D}\right].
\]
Выбираем ветвь \(\sqrt D=1\) при \(t=0\).
Для потока положим \(u=x-x_0,\ v=x+x_0,\ B=1+\tau^2\) и
\[
 R=e^{-(x^2+x_0^2)/(\sigma^2 B)},\qquad
 \theta=\frac{2\tau x x_0}{\sigma^2 B}.
\]
Тогда \(g_+^*g_-=Re^{i\theta}\).
Дифференцирование даёт \(\partial_x\Psi
=-C(u g_++v g_-)/(\sigma^2 D^{3/2})\).
В произведении \((g_+^*+g_-^*)(u g_++v g_-)\)
вещественная часть равна
\(u e^{-u^2/(\sigma^2 B)}+v e^{-v^2/(\sigma^2 B)}+2xR\cos\theta\),
а мнимая равна \(2x_0 R\sin\theta\). Поэтому
\begin{align*}
 j(x,t)&=\frac{\hbar C^2}{m\sigma^2 B^{3/2}}\Big[
 \tau\big(u e^{-u^2/(\sigma^2B)}+v e^{-v^2/(\sigma^2B)}
                   +2xR\cos\theta\big)-2x_0R\sin\theta\Big].
\end{align*}
Поток равен нулю при \(t=0\) и нечётен по \(x\) при всех \(t\),
что согласуется с чётностью состояния.

\end{document}
